\documentclass[11pt]{article}
\usepackage[margin=1in]{geometry}
\usepackage[T1]{fontenc}
\usepackage[utf8]{inputenc}
\usepackage{lmodern}
\usepackage{microtype}
\usepackage{xcolor}
\usepackage{amsmath,amssymb,amsthm,mathtools}
\usepackage{bm}
\usepackage{hyperref}
\usepackage{enumitem}
\usepackage{booktabs}
\usepackage{array}
\usepackage{listings}
\usepackage{inconsolata}
\usepackage{fancyhdr}
\usepackage{titlesec}

\hypersetup{
  colorlinks=true,
  linkcolor=blue!50!black,
  urlcolor=blue!50!black,
  citecolor=blue!50!black,
  pdftitle={Middle Product via FFT and MSM},
  pdfauthor={OpenAI}
}

\pagestyle{fancy}
\fancyhf{}
\lhead{Middle Product via FFT and MSM}
\rhead{\thepage}
\setlength{\headheight}{14pt}

\titleformat{\section}{\large\bfseries\color{blue!35!black}}{\thesection.}{0.5em}{}
\titleformat{\subsection}{\normalsize\bfseries\color{blue!20!black}}{\thesubsection.}{0.5em}{}

\lstdefinestyle{mystyle}{
  basicstyle=\ttfamily\small,
  keywordstyle=\bfseries\color{blue!50!black},
  commentstyle=\itshape\color{black!55},
  frame=single,
  framerule=0.4pt,
  rulecolor=\color{black!20},
  backgroundcolor=\color{black!2},
  showstringspaces=false,
  columns=fullflexible,
  keepspaces=true,
  xleftmargin=0.5em,
  xrightmargin=0.5em
}
\lstset{style=mystyle}

\newcommand{\GT}{G_T}
\newcommand{\cA}{\mathcal{A}}
\newcommand{\cB}{\mathcal{B}}
\newcommand{\cC}{\mathcal{C}}

\begin{document}

\begin{center}
  {\LARGE \textbf{Middle Product via FFT, Pairings, and MSM}}\\[0.5em]
  {\large Technical note}\\[0.75em]
\end{center}

\noindent
This note explains the convolution view, the optimized middle-product FFT view, the direct quadratic alternative, and a practical FFT/MSM hybrid. The intended audience already understands the batching use case and the standard FFT-based acceleration, so the emphasis here is on implementable detail.

\section{Problem statement}

Let
\[
A(x) = a_0 + a_1 x + \cdots + a_{n-1} x^{n-1}
\]
and
\[
B(x) = b_0 + b_1 x + \cdots + b_{2n-1} x^{2n-1},
\]
where the coefficients of $A$ live in one pairing input group, the coefficients of $B$ live in the other pairing input group, and coefficient multiplication is the pairing
\[
a_i \circ b_j \in \GT.
\]

Define
\[
C(x) = A(x)B(x) = \sum_{k=0}^{3n-2} C_k x^k.
\]
We do \emph{not} need the full product. We only need the middle window
\[
C_n, C_{n+1}, \dots, C_{2n-1}.
\]

This is exactly the middle-product problem.

\section{Naive formula}

For each $r \in \{0,\dots,n-1\}$,
\[
C_{n+r} = \sum_{i=0}^{n-1} a_i \circ b_{n+r-i}.
\]
Because $B$ has length $2n$, these are full length-$n$ sums with no boundary truncation in the middle window.

The direct implementation is therefore:

\begin{lstlisting}
for r in 0..n-1:
    acc = 0 in G_T
    for i in 0..n-1:
        acc += pair(a[i], b[n + r - i])
    out[r] = acc
\end{lstlisting}

This costs $n^2$ bilinear products. It is the cleanest statement of what the middle product is computing.

\section{Basic FFT version: compute the whole product}

The standard convolution identity still applies:
\begin{enumerate}[itemsep=0.3em, topsep=0.4em]
  \item transform $A$ and $B$ to the evaluation domain,
  \item multiply pointwise there,
  \item inverse transform back.
\end{enumerate}

Since
\[
\deg(A) \le n-1,
\qquad
\deg(B) \le 2n-1,
\]
we have
\[
\deg(C) \le 3n-2,
\]
so the full linear convolution has length
\[
3n-1.
\]
Therefore the basic FFT implementation chooses a transform length
\[
m_{\mathrm{basic}} \ge 3n-1,
\]
typically the next power of two, often about $4n$.

Let $m = m_{\mathrm{basic}}$, let $\omega$ be a primitive $m$th root of unity in the scalar field, and zero-pad $A$ and $B$ to length $m$. Then compute
\[
\widehat{A} = \mathrm{FFT}_m(A),
\qquad
\widehat{B} = \mathrm{FFT}_m(B),
\]
then pointwise pair/multiply
\[
\widehat{C}_j = \widehat{A}_j \circ \widehat{B}_j,
\qquad j = 0,\dots,m-1,
\]
and finally recover the coefficient vector
\[
C = \mathrm{iFFT}_m(\widehat{C}).
\]

The online flow is:
\begin{enumerate}[itemsep=0.3em, topsep=0.4em]
  \item pad $A$ and $B$ to size $m \ge 3n-1$,
  \item FFT $A$,
  \item FFT $B$,
  \item pointwise pair/multiply,
  \item inverse FFT,
  \item discard everything except coefficients $[n,2n)$.
\end{enumerate}

This works, but it pays for the full product.

\section{Optimized middle-product version}

The key improvement is not merely ``precompute $B$.'' The key improvement is that the middle-product algebra lets us \emph{shrink the transform length}.

\subsection{Split $B$ into two halves}

Write
\[
B(x) = B_0(x) + x^n B_1(x),
\]
where
\[
B_0(x) = b_0 + b_1 x + \cdots + b_{n-1} x^{n-1},
\]
and
\[
B_1(x) = b_n + b_{n+1} x + \cdots + b_{2n-1} x^{n-1}.
\]
Then
\[
A(x)B(x) = A(x)B_0(x) + x^n A(x)B_1(x).
\]

Now observe:
\begin{align*}
\deg(A B_0) &\le 2n-2,\\
\deg(A B_1) &\le 2n-2.
\end{align*}
So after multiplying the second term by $x^n$, it lives in degrees
\[
[n, 3n-2].
\]
We only care about degrees
\[
[n, 2n-1].
\]

\subsection{Why size $m \ge 2n$ is enough}

Work modulo $x^m - 1$ for any
\[
m \ge 2n.
\]
Then:
\begin{itemize}[itemsep=0.3em, topsep=0.4em]
  \item $A(x)B_0(x)$ does not wrap at all, because its degree is $< 2n \le m$.
  \item $x^n A(x)B_1(x)$ may wrap, but only terms of degree $\ge m$ wrap.
  \item After reduction modulo $x^m-1$, those wrapped terms land in degrees $< n$.
\end{itemize}

Therefore any aliasing caused by using an $m$-point \emph{cyclic} convolution can only pollute coefficients
\[
0,1,\dots,n-1,
\]
and cannot affect the middle window
\[
n,n+1,\dots,2n-1.
\]

Hence:

\begin{quote}
The middle window of the true product agrees exactly with the middle window of the cyclic product modulo $x^m-1$, for any $m \ge 2n$.
\end{quote}

This is the central algebraic fact behind the optimized version.

So the optimized implementation only needs transform length
\[
m_{\mathrm{opt}} \ge 2n,
\]
typically the next power of two, often about $2n$ instead of $4n$.

\subsection{Evaluation-domain combination}

Let $w_j$ be the $j$th FFT evaluation point for the chosen length $m$. At each evaluation point,
\[
B(w_j) = B_0(w_j) + w_j^n B_1(w_j).
\]
So after precomputing the FFTs of the two halves,
\[
\widehat{B_0}_j = B_0(w_j),
\qquad
\widehat{B_1}_j = B_1(w_j),
\]
we combine them online as
\[
\widehat{B}_j = \widehat{B_0}_j + w_j^n \widehat{B_1}_j.
\]

This is the ``combine the halves in the evaluation domain'' step.

\subsection{Online path}

The optimized online path is:
\begin{enumerate}[itemsep=0.3em, topsep=0.4em]
  \item choose $m = \mathrm{pow2ceil}(2n)$,
  \item FFT $A$ at size $m$,
  \item combine the two precomputed halves of $B$:
  \[
  \widehat{B}_j = \widehat{B_0}_j + w_j^n \widehat{B_1}_j,
  \]
  \item pointwise pair:
  \[
  \widehat{C}_j = \widehat{A}_j \circ \widehat{B}_j,
  \]
  \item inverse FFT,
  \item keep only coefficients $[n,2n)$.
\end{enumerate}

So compared with the basic full-product path, the optimized path gets:
\begin{itemize}[itemsep=0.3em, topsep=0.4em]
  \item smaller FFT size,
  \item fewer pointwise pairings,
  \item a smaller inverse FFT,
  \item and, if desired, a truncated inverse FFT that recovers only the middle window.
\end{itemize}

\section{Concrete implementation structure}

\subsection{Offline precomputation for $B$}

Choose
\[
m = \mathrm{pow2ceil}(2n).
\]
Build the two halves $B_0, B_1$, pad each to length $m$, and compute
\[
\widehat{B_0} = \mathrm{FFT}_m(B_0),
\qquad
\widehat{B_1} = \mathrm{FFT}_m(B_1).
\]
Also precompute the twist vector
\[
u_j = w_j^n,
\qquad j=0,\dots,m-1.
\]
Store $(m, \widehat{B_0}, \widehat{B_1}, u)$.

\subsection{Online computation for each $A$}

Given $A$, zero-pad it to length $m$ and compute
\[
\widehat{A} = \mathrm{FFT}_m(A).
\]
Then for each $j$:
\[
\widehat{B}_j = \widehat{B_0}_j + u_j \widehat{B_1}_j,
\qquad
\widehat{C}_j = \widehat{A}_j \circ \widehat{B}_j.
\]
Then invert:
\[
C = \mathrm{iFFT}_m(\widehat{C}),
\]
and return
\[
(C_n,\dots,C_{2n-1}).
\]

Pseudocode:

\begin{lstlisting}
offline_precompute_B(B):
    split B into B0, B1, each length n
    m = next_power_of_two(2*n)

    B0_pad = zero_pad(B0, m)
    B1_pad = zero_pad(B1, m)

    B0_fft = fft_group(B0_pad)
    B1_fft = fft_group(B1_pad)

    roots = fft_roots(m)
    twist[j] = roots[j]^n for j in 0..m-1

    return (m, B0_fft, B1_fft, twist)
\end{lstlisting}

\begin{lstlisting}
middle_product_online(A, precomp):
    (m, B0_fft, B1_fft, twist) = precomp

    A_pad = zero_pad(A, m)
    A_fft = fft_group(A_pad)

    for j in 0..m-1:
        B_fft[j] = B0_fft[j] + twist[j] * B1_fft[j]
        C_fft[j] = pair(A_fft[j], B_fft[j])

    C = ifft_target_group(C_fft)

    return C[n .. 2*n-1]
\end{lstlisting}

\section{How FFT stages become MSMs}

An FFT stage is just a fixed linear transform over group-valued data.

For radix $2$, the butterfly is
\[
x' = x + \omega y,
\qquad
y' = x - \omega y.
\]
When $x$ and $y$ are group elements, these are scalar-group linear combinations:
\[
x' = 1 \cdot x + \omega \cdot y,
\qquad
y' = 1 \cdot x - \omega \cdot y.
\]
So each output is already a tiny multiscalar multiplication.

More generally, a radix-$r$ stage applies a fixed $r \times r$ matrix $M_r$:
\[
y_u = \sum_{v=0}^{r-1} M_r[u,v] \, x_v.
\]
Each output $y_u$ is an MSM over the $r$ inputs.

So there are two equivalent implementation views:
\begin{itemize}[itemsep=0.3em, topsep=0.4em]
  \item use explicit butterflies,
  \item or view each stage as a fixed linear map and evaluate it with MSM kernels.
\end{itemize}

The algebra is identical. FFT provides the structured linear map; MSM is one possible primitive for evaluating it.

\section{Direct quadratic approach via multipairing / MSM-style batching}

The naive formula
\[
C_{n+r} = \sum_{i=0}^{n-1} a_i \circ b_{n+r-i}
\]
can also be implemented as a dense batched kernel.

For each output $r$, form the $n$ pairs
\[
(a_0,b_{n+r}),\ (a_1,b_{n+r-1}),\ \dots,\ (a_{n-1},b_r),
\]
and compute one multipairing:
\[
C_{n+r} = \mathrm{MultiPair}\big((a_i,b_{n+r-i})_{i=0}^{n-1}\big).
\]

That is still $O(n^2)$ total work, but much cheaper than doing $n^2$ completely separate pairings.

This is why the direct method is not crazy for small $n$. Multipairing and MSM-style kernels are often highly optimized. For small batches, constant factors can beat the asymptotically better FFT path.

A reasonable engineering strategy is:
\begin{itemize}[itemsep=0.3em, topsep=0.4em]
  \item benchmark direct multipairing for small $n$,
  \item switch to the FFT path once $n$ crosses the empirical threshold.
\end{itemize}

\section{FFT/MSM hybrid}

This naturally suggests a hybrid implementation.

At large sizes, use the FFT structure, since that gives the correct asymptotic behavior. At small sizes, do not recurse all the way down to tiny radix-$2$ butterflies. Instead, stop at some cutoff $r \in \{8,16,32\}$ and replace the remaining local transform by a dense radix-$r$ kernel implemented via MSM.

Equivalently:
\begin{itemize}[itemsep=0.3em, topsep=0.4em]
  \item collapse the last $\log_2 r$ radix-$2$ layers into one radix-$r$ stage,
  \item compute that stage output-by-output as small MSMs.
\end{itemize}

This can be viewed as a \emph{high-radix FFT} whose butterfly logic is implemented by MSM kernels rather than explicit hand-written circuits.

So the clean conceptual picture is:
\begin{enumerate}[itemsep=0.3em, topsep=0.4em]
  \item overall flow remains
  \[
  \text{convolution on } A,B \;\to\; \text{pairing in the evaluation domain} \;\to\; \text{inverse convolution on } C,
  \]
  \item the optimized middle-product algebra shrinks the transform length from about $4n$ to about $2n$,
  \item the direct $O(n^2)$ path is still viable for small $n$ because multipairing and MSM are fast,
  \item and the most practical implementation may be a hybrid FFT/MSM scheme, viewed as a high-radix FFT with cleverly MSM-implemented leaf stages.
\end{enumerate}

\section{Short intuition}

The optimized trick works because the upper half of $B$ appears with an $x^n$ shift:
\[
B(x) = B_0(x) + x^n B_1(x).
\]
When we compress the computation into a size-$m$ cyclic domain with $m \ge 2n$, any aliasing from that shifted high part wraps only into degrees $< n$. Since we only read degrees $[n,2n)$, that aliasing is harmless.

That is why the middle product can be computed with transform size about $2n$ instead of full-product size about $4n$.

\section{Implementation note.}
The FFT/MSM hybrid does not replace the middle-product algorithm itself.
It only replaces selected group-FFT stages.
At a cutoff radix $r$, we stop expanding radix-$2$ butterflies and instead
evaluate the remaining local size-$r$ transform as a dense fixed linear map
over group elements.

If the stage input is $(x_0,\dots,x_{r-1})$ and the local transform matrix is
$M \in \mathbb{F}_r^{r \times r}$, then each output
\[
y_u = \sum_{v=0}^{r-1} M[u,v] x_v
\]
is an MSM over the shared base set $(x_0,\dots,x_{r-1})$.
Thus a radix-$r$ leaf is naturally implemented as a batched family of MSMs
with common bases and fixed coefficients.
In the limit $r=2$, this reduces to the usual radix-$2$ butterfly network;
in the limit $r=m$, it becomes a dense full-transform implementation.

\section{Implementation plan}

The engineering goal is to build a \emph{generalized group-transform kernel}
first, and only then plug it into the middle-product path.
We are \emph{not} implementing a new middle product.
We are implementing an alternative way to evaluate the same group FFT.

\subsection{Phase 1: a standalone group-linear kernel}

The first deliverable is a routine that takes $m$ group inputs and produces
$m$ group outputs:
\[
(x_0,\dots,x_{m-1}) \mapsto (y_0,\dots,y_{m-1}),
\qquad
y_u = \sum_{v=0}^{m-1} M[u,v] x_v.
\]
Here $M$ is a fixed transform matrix coming from a dense FFT leaf or from the
full transform itself.

The important interface point is that we do not need to pass arbitrary field
elements at runtime.
For the planned use case, the coefficient pattern is fixed by the transform.
So the kernel can accept coefficient \emph{indexes} or twiddle-table
references instead of an explicit dense matrix of field elements.

A practical API shape is:
\begin{lstlisting}
group_linear_transform(
    out[0..m-1],
    in[0..m-1],
    transform_id,
    twiddle_index_table)
\end{lstlisting}
where the transform descriptor identifies which fixed linear map is being
applied and the index table tells each output row which precomputed scalar to
use for each input.

Conceptually, each row is an MSM:
\[
y_u = \sum_{v=0}^{m-1} \lambda_{u,v} x_v.
\]
Operationally, the implementation should treat the whole stage as a
\emph{batched shared-base} problem, since all $m$ outputs use the same base set
$(x_0,\dots,x_{m-1})$.

\subsection{Phase 2: correctness against the existing FFT}

Before touching the middle-product code, benchmark and test the standalone
kernel against the existing radix-$2$ group FFT.

For each target size $m$:
\begin{enumerate}[itemsep=0.3em, topsep=0.4em]
  \item generate deterministic random group inputs,
  \item compute the reference result with the current FFT implementation,
  \item compute the same transform with the new group-linear kernel,
  \item compare outputs elementwise.
\end{enumerate}

This isolates the new building block.
If it is wrong or slow, we learn that before any middle-product integration
complicates the picture.

\subsection{Phase 3: determine the crossover range}

Once correctness is in place, the next task is pure benchmarking.
Measure the standalone group-linear kernel against the existing radix-$2$ FFT
for a sweep of sizes such as
\[
m \in \{2,4,8,16,32,64,128\}.
\]

The main question is:
\begin{quote}
For which transform sizes is the dense shared-base MSM-style kernel faster
than the radix-$2$ FFT implementation?
\end{quote}

This experiment should be run on real AVX hardware.
The output of this phase is not a new algorithm; it is a cutoff table or at
least an empirical threshold.

\subsection{Phase 4: use the data to choose a hybrid cutoff}

The benchmark data from Phase 3 should determine where the hybrid begins and
ends.
For example, if the dense kernel wins up to $m=16$ and loses at $m=32$, then a
natural first hybrid is:
\begin{itemize}[itemsep=0.3em, topsep=0.4em]
  \item keep radix-$2$ FFT layers above size $16$,
  \item replace the final four radix-$2$ layers by one dense size-$16$ leaf.
\end{itemize}

That gives a concrete high-radix leaf policy derived from measurements rather
than from guesswork.

\subsection{Phase 5: integrate into the optimized middle-product path}

Only after the leaf-kernel crossover is understood should the hybrid be wired
into the middle-product implementation.
At that point, the middle-product structure remains unchanged:
\begin{enumerate}[itemsep=0.3em, topsep=0.4em]
  \item use the optimized size-$2n$ middle-product algebra,
  \item transform $A$ with the hybrid group FFT,
  \item evaluate the precomputed $B$ side in the same transform basis,
  \item pair pointwise in the evaluation domain,
  \item run the inverse hybrid transform or truncated inverse as before.
\end{enumerate}

So the integration step is deliberately late in the plan.
The purpose of the first phases is to avoid mixing up two independent
questions:
\begin{itemize}[itemsep=0.3em, topsep=0.4em]
  \item whether the dense shared-base group-linear kernel is actually good,
  \item and how best to use it inside the middle-product pipeline.
\end{itemize}

\subsection{Phase 6: optional endpoint experiments}

Two endpoint experiments are useful for calibration:
\begin{itemize}[itemsep=0.3em, topsep=0.4em]
  \item \textbf{Pure radix-$2$ FFT:} the current baseline.
  \item \textbf{Pure dense transform:} compute the whole size-$m$ transform with
  the new group-linear kernel and no radix-$2$ recursion.
\end{itemize}

These endpoints are valuable even if only the hybrid is ultimately used,
because they show whether the interpolation curve is smooth and whether the
chosen leaf cutoffs are close to optimal.


\end{document}
